\chapter{Discussion}

Throughout this work, we have investigated how different assumptions regarding the high-energy behaviour of scattering amplitudes constrain the low-energy spectrum in various theoretical settings.

First, we demonstrated that a theory of massless scalars satisfying $-2$-subtracted dispersion relations cannot feature the exchange of a massive resonance with any spin. For scalar fields of mass $M$ obeying the same dispersion relations, the exchange of an even-spin resonance of mass $m$ is permitted only if
\[
m^2 \geq \frac{M^2}{2}.
\]

We then extended the analysis to theories satisfying $0$-subtracted dispersion relations and including the exchange of spin-$0$ and spin-$2$ resonances. Analytically, we obtained the inequality
\[
m_0 < m_2,
\]
and found that the exchange of a spin-$2$ resonance is impossible in the absence of a scalar resonance at a lower energy. Moreover, our numerical analysis provided upper bounds on the coupling of the spin-$2$ resonance to the external scalar states and reinforced the conclusion that a spin-$0$ resonance is necessary for a consistent heavier spin-$2$ state. The numerical bounds also impose stronger constraints on the allowed mass spectrum than those derived analytically. Similar results were obtained for the case of massive scalar external legs.

Then, we explored how to obtain a general parametrisation of the scattering amplitude for non-abelian massive vectors. After understanding the role of crossing symmetry in its construction, we described a simplified method for deriving positivity bounds both in the forward limit and for finite values of $t$.

We applied these results to study $WW$ scattering. First, we found that no deformation of the couplings and masses of the electroweak sector of the Standard Model is compatible with $0$-subtracted dispersion relations. However, compatibility can be restored by fine-tuning the contributions from higher-order operators of Beyond the Standard Model physics. Nevertheless, within the bounds we employed, we were unable to derive bounds on the spectrum that were independent of the operators. We also examined the case in which the process satisfies once-subtracted dispersion relations, assuming analyticity in the isospin-$2$ channel. The Standard Model was found to be compatible with this setup, and we determined the allowed parameter space for the contributions from higher-order operators.